\documentclass[11pt,a4paper]{article}
\usepackage{fontspec}
\setmainfont{Noto Sans SC}[
  Path=fonts/package/,
  UprightFont=400Regular/NotoSansSC_400Regular.ttf,
  BoldFont=700Bold/NotoSansSC_700Bold.ttf
]
\usepackage[margin=2.35cm]{geometry}
\usepackage{amsmath,amssymb,amsthm,mathtools,bm}
\usepackage{enumitem,booktabs,array,longtable,tabularx}
\usepackage{xcolor,hyperref,cleveref,fancyhdr,titlesec}
\definecolor{deepblue}{HTML}{16324F}
\definecolor{softblue}{HTML}{EAF2F8}
\definecolor{warm}{HTML}{8A4B08}
\hypersetup{colorlinks=true,linkcolor=deepblue,citecolor=warm,urlcolor=deepblue,
  pdftitle={Analysis on the Levi--Civita Field: Reading Notes},
  pdfauthor={Personal reading notes}}
\setlength{\parindent}{1.5em}
\setlength{\parskip}{0.35em}
\setlength{\emergencystretch}{3em}
\setlist{nosep,leftmargin=2.2em}
\sloppy
\pagestyle{fancy}
\setlength{\headheight}{14pt}
\fancyhf{}
\fancyhead[L]{Analysis on the Levi--Civita Field}
\fancyhead[R]{Personal Reading Notes}
\fancyfoot[C]{\thepage}

\newtheorem{definition}{Definition}[section]
\newtheorem{theorem}[definition]{Theorem}
\newtheorem{proposition}[definition]{Proposition}
\newtheorem{lemma}[definition]{Lemma}
\newtheorem{corollary}[definition]{Corollary}
\theoremstyle{remark}
\newtheorem{remark}[definition]{Remark}
\newtheorem{example}[definition]{Example}

\newcommand{\LC}{\mathcal R}
\newcommand{\supp}{\operatorname{supp}}
\newcommand{\ord}{\operatorname{ord}}
\newcommand{\Aut}{\operatorname{Aut}}
\newcommand{\R}{\mathbb R}
\newcommand{\Q}{\mathbb Q}
\newcommand{\Z}{\mathbb Z}
\newcommand{\C}{\mathbb C}
\newcommand{\st}{\operatorname{st}}

\title{\bfseries\color{deepblue} Analysis on the Levi--Civita Field\par
\large From Genuine Infinitesimals and Generalized Series to Computable Differentiation\par
\vspace{0.4em}\normalsize A Guided Reading of Khodr M. Shamseddine's Dissertation\par
\textit{New Elements of Analysis on the Levi-Civita Field}}
\author{Personal Reading Notes}
\date{July 2026}

\begin{document}
\maketitle

\begin{center}
\fcolorbox{deepblue}{softblue}{\parbox{0.88\textwidth}{
\textbf{Authorship and provenance.} These are personal study notes based on a publicly
available dissertation and earlier learning discussions. They are not the note writer's
research and do not imply participation in the original work. The dissertation was first
recommended by a friend's advisor to that friend; the note writer pursued it independently
out of personal interest. The original work is Khodr M. Shamseddine's 1999 Michigan State
University doctoral dissertation, supervised by Martin Berz. The author's copy is available
at \href{https://www.physics.umanitoba.ca/~khodr/dissert.pdf}{dissert.pdf}.
}}
\end{center}

\begin{abstract}
These notes organize the central questions raised by Shamseddine's dissertation into one
coherent route. Why leave the Archimedean real field? What is gained by admitting genuine
infinitesimals? How is the Levi--Civita field constructed, and why is left-finite support the
right compromise between Hahn-series generality and effective computation? What do leading
exponents, valuations, skeleton groups, field automorphisms, real closure, differential
algebra, order topology, and weak topology contribute? The later sections explain why naive
topological continuity and differentiability are too weak, why uniform Lipschitz-type control
and expandable functions are introduced, and how computer functions and infinitesimal
differentiation connect the abstract theory to high-order numerical work. Common
misconceptions about point particles, Dirac delta, finite support, and weak convergence are
addressed explicitly.
\end{abstract}

\tableofcontents
\newpage

\section{The global problem}

The dissertation does not merely replace one derivative formula by another. It builds an
environment in which infinitesimals are numbers, generalized series remain computable, and a
useful part of classical analysis can be reconstructed. Its logical architecture is
\[
\begin{gathered}
\boxed{\text{genuine infinitesimals}}
\Longrightarrow
\boxed{\text{a computable non-Archimedean field}}\\[0.4em]
\Downarrow\\[-0.1em]
\boxed{\text{calculus and applications}}
\Longleftarrow
\boxed{\text{appropriate topologies and function classes}}.
\end{gathered}
\]
The point is not that real analysis is defective. Real analysis is rigorous and extraordinarily
successful. The new theory answers a different design question: can infinitesimal reasoning be
made literal while retaining algebraic control, convergence mechanisms, and implementation?

The chapters play distinct roles. Chapter 1 supplies mathematical, physical, and computational
motivation. Chapter 2 studies the algebraic skeleton of non-Archimedean extensions. Chapter 3
constructs the Levi--Civita field and its order, topology, and differential algebra. Chapter 4
develops sequences and power series. Chapter 5 shows why naive calculus is pathological.
Chapter 6 isolates expandable functions, a class on which familiar theorems return. Chapter 7
connects the theory with computer functions and higher-order differentiation.

\section{Why leave the real field? Three motivations}

\subsection{The mathematical motivation}

\begin{definition}[Archimedean ordered field]
An ordered field $F$ is Archimedean if for every $x\in F$ there exists $n\in\mathbb N$ with
$n>x$. Equivalently, no positive $\varepsilon\in F$ satisfies
\[
0<\varepsilon<\frac1n\qquad\text{for every }n\in\mathbb N.
\]
\end{definition}

The real field is Archimedean. Consequently a nonzero number smaller than every positive real
scale cannot belong to $\R$. Classical analysis handles this by limits:
\[
f'(x)=\lim_{h\to0}\frac{f(x+h)-f(x)}{h}.
\]
The dissertation asks whether one can work in a larger ordered field containing a nonzero $d$
with $|d|<1/n$ for every $n$, and then manipulate the difference quotient directly. This is an
extension of the number system, not a criticism of the limit definition.

\subsection{The computational motivation}

Ordinary finite differences face two competing errors. If the real step $h$ is too large, Taylor
truncation dominates; if it is too small, subtracting nearly equal floating-point values causes
catastrophic cancellation. With a formal genuine infinitesimal $d$, different error orders are
stored in different powers of $d$. For example,
\[
f(x)=x^2-2x
\quad\Longrightarrow\quad
\frac{f(x+d)-f(x)}d=2x-2+d.
\]
The quotient is not literally the derivative: it is the derivative plus an infinitesimal
remainder. Reading its constant coefficient gives $2x-2$. This coefficient-extraction viewpoint
is the bridge from infinitesimal algebra to automatic or numerical differentiation.

\subsection{The physical motivation}

Point charges, point masses, and impulses are represented in classical analysis by distributions.
The Dirac delta is defined by its action on a test function,
\[
\langle\delta_0,\varphi\rangle=\varphi(0),
\]
not as an ordinary real-valued function. A non-Archimedean framework can instead use an actual
function with infinitesimal width and infinite height whose integral action reproduces the desired
point-source behavior.

This does not mean that an electron or a smallest physical unit is a mathematical infinitesimal.
An elementary charge is a finite physical constant; a point particle is a geometric idealization;
a Levi--Civita infinitesimal is a nonzero element of an ordered field. These are three different
logical categories. The connection is modeling: infinitesimal-width profiles may represent
concentrated sources while remaining genuine functions in the enlarged codomain.

\section{From design requirements to the field}

Before writing a formula, list the desired properties:
\begin{enumerate}
\item $\R$ should embed as an ordered subfield.
\item Positive infinitesimals $d$ and rational powers $d^q$ should exist.
\item Every nonzero element should have a readable dominant term.
\item Addition and multiplication should be computable coefficient by coefficient.
\item Each output coefficient should depend on only finitely many input coefficients.
\item The resulting structure should be a field and, ideally, real closed and suitably complete.
\end{enumerate}

These requirements suggest generalized series
\[
x=\sum_{q\in\Q}x[q]d^q,
\]
where $x[q]\in\R$. The crucial issue is not writing the symbol but controlling the support
\[
\supp(x)=\{q\in\Q:x[q]\ne0\}.
\]

\begin{definition}[Left-finite support]
A set $S\subseteq\Q$ is left-finite if, for every $r\in\Q$, the set
\[
S\cap(-\infty,r]
\]
is finite.
\end{definition}

Left-finite does \emph{not} mean finite. The support $\{0,1,2,\ldots\}$ is infinite but
left-finite. In contrast, $\{1,1/2,1/3,\ldots\}$ is not left-finite because infinitely many
terms lie to the left of $1$. The adjective ``left'' refers to the order direction on the
exponent line, not an arbitrary naming convention.

\begin{definition}[Levi--Civita field]
The real Levi--Civita field is
\[
\LC=\{x:\Q\to\R\mid\supp(x)\text{ is left-finite}\}.
\]
It is convenient to read $x$ as the generalized series $\sum_q x[q]d^q$.
\end{definition}

Addition is coefficientwise. Multiplication is convolution:
\[
(xy)[q]=\sum_{q_1+q_2=q}x[q_1]y[q_2].
\]
The left-finite condition ensures that, for fixed $q$, only finitely many pairs contribute.
Thus each coefficient is an ordinary finite real sum. This is the computational reason for the
support restriction.

Hahn fields permit much more general well-ordered supports. Every left-finite subset of $\Q$ is
well ordered, but not every well-ordered subset is left-finite. Hence the Levi--Civita field is a
special Hahn-type generalized-series field. It sacrifices some generality to obtain sequential,
coefficientwise computability.

\section{Leading exponents, order, and valuation}

For nonzero $x\in\LC$, define
\[
\lambda(x)=\min\supp(x).
\]
This leading exponent gives the dominant scale. If
\[
x=7d^{3/2}-2d^2+5d^{11/4},
\]
then $3/2$ has not appeared mysteriously: it is simply the smallest exponent with nonzero
coefficient, so $\lambda(x)=3/2$. The element is of order $d^{3/2}$, and
\[
d^{-3/2}x=7-2d^{1/2}+5d^{5/4}
\]
has finite nonzero leading part $7$.

The field order is lexicographic by the leading coefficient: $x>0$ if and only if
$x[\lambda(x)]>0$. This works because all higher-exponent terms are smaller scales and cannot
reverse the sign of the first nonzero term.

The map $v(x)=\lambda(x)$ behaves like a valuation:
\[
v(xy)=v(x)+v(y),\qquad
v(x+y)\ge\min\{v(x),v(y)\}.
\]
It is best understood as order-of-magnitude bookkeeping. Large valuation means small magnitude.
For instance, $v(d^5)>v(d^2)$ but $0<d^5<d^2$.

Two positive elements are Archimedean equivalent when each is bounded by a natural-number
multiple of the other. The equivalence classes form a group of magnitude levels. A
\emph{skeleton group} is a structural record of these levels together with the coefficient data
attached to them. It separates ``which scale?'' from ``which leading coefficient?'' and helps
classify automorphisms and extensions.

\section{Algebraic background with minimal prerequisites}

A field is a commutative number system in which every nonzero element has a multiplicative
inverse. A field automorphism is a bijection $\sigma:F\to F$ preserving addition, multiplication,
$0$, and $1$. Every field automorphism fixes $\Q$, because it fixes $1$ and therefore all integer
and rational combinations of $1$.

``Nontrivial automorphism'' means an automorphism other than the identity. The phrase does not
mean that the field itself is nontrivial. When an order must also be preserved, the freedom is
smaller: positive elements must stay positive and the hierarchy of scales must remain compatible.
In a non-Archimedean field, understanding possible changes of infinitesimal generators requires
understanding the value group and its skeleton.

An algebraically closed field is one in which every nonconstant polynomial has a root. The real
numbers are not algebraically closed because $X^2+1$ has no real root; the complex numbers are.
For ordered fields the more natural endpoint is \emph{real closed}: every positive element has a
square root and every odd-degree polynomial has a root; adjoining $i$ then gives an algebraically
closed field. Real closure matters because it ensures that algebraic operations do not continually
force us outside the chosen real domain.

A differential field is a field equipped with a derivation $D$ satisfying
\[
D(x+y)=D(x)+D(y),\qquad D(xy)=xD(y)+yD(x).
\]
For generalized series one can formally set $D(d^q)=q d^{q-1}$ and extend termwise under the
support rules. This is a differential-algebra structure on the numbers themselves. It is distinct
from differentiating a function $f:\LC\to\LC$, although the two ideas later interact.

\section{Order topology and weak topology}

The order topology is generated by open intervals $(a,b)$. Because the field contains genuine
infinitesimals, it is much finer near real points than the usual topology of $\R$. In particular,
the embedded copy of $\R$ is discrete: for a positive infinitesimal $d$,
\[
(r-d,r+d)\cap\R=\{r\}.
\]
This explains why connectedness-based real-analysis arguments cannot be transferred blindly.
An interval in a non-Archimedean ordered field may be disconnected in ways impossible over $\R$,
and local compactness may fail because infinitesimal scales create too many nested directions.

Order-Cauchy completeness means that every sequence Cauchy with respect to the ordered-field
absolute value converges in the field. It is useful, but it is not the same as Dedekind
completeness. The latter would force an Archimedean ordered field under familiar hypotheses, so
one must specify the completeness notion rather than say simply ``complete.''

The weak topology observes finitely many coefficients at a time. One useful family of seminorms
is
\[
p_r(x)=\max\{|x[q]|:q\le r\},
\]
with the convention that the maximum over an empty set is zero. Left-finiteness makes the set of
relevant nonzero coefficients finite, so the maximum exists. The seminorm axioms follow directly:
\[
p_r(\alpha x)=|\alpha|p_r(x),\qquad
p_r(x+y)\le p_r(x)+p_r(y),\qquad p_r(x)\ge0.
\]
It may vanish on a nonzero element: if all exponents of $x$ exceed $r$, then $p_r(x)=0$. This is
exactly why it is a seminorm, not a norm.

Weak convergence $x_n\to x$ means convergence with respect to all these seminorms, equivalently
coefficientwise convergence on every bounded-left exponent window. Strong or order convergence
controls the whole ordered magnitude. The philosophical resemblance to weak topologies in
functional analysis is real: both retain selected observables while permitting convergence that
a stronger topology would reject. The concrete observables, however, are different.

\section{Why ordinary continuity and differentiability become pathological}

Classical conclusions often contain hidden topological hypotheses. On a connected real interval,
a differentiable function with zero derivative is constant. More generally, on a connected open
subset of $\R^n$, a $C^1$ function with zero differential is constant. In complex analysis, a
holomorphic function with zero complex derivative is constant on each connected component. The
proof always uses a mechanism linking nearby points---usually the mean value theorem, paths, or
analytic continuation.

Over a highly disconnected non-Archimedean domain, a locally constant function may vary from one
clopen component to another while having derivative zero everywhere. Thus the formal statement
``$f'=0$ implies $f$ is constant'' fails unless the domain and function class restore enough
coherence. Likewise, compactness-based results such as uniform continuity on closed bounded
intervals, attainment of extrema, and subsequence extraction may fail when the corresponding
topological compactness is absent.

Topological continuity alone only says that each point has some sufficiently small neighborhood.
The permitted scale can vary wildly from point to point. A uniform Lipschitz condition
\[
|f(x)-f(y)|\le L|x-y|
\]
uses one constant $L$ across the relevant set and couples behavior at different points. It is not
an arbitrary strengthening: it supplies exactly the uniform control needed for estimates,
integration, limit interchange, and mean-value-type arguments.

This design pattern occurs throughout analysis. In PDE one rarely works with all functions:
$C^k$ spaces control classical derivatives, $C^\infty$ supports repeated differentiation,
analytic functions are controlled by convergent power series, Sobolev spaces control weak
derivatives in $L^p$, and H\"older or Lipschitz spaces control oscillation. A ``good function
class'' is chosen to be strong enough for theorems and stable under the operations the problem
requires.

\section{Derivate and expandable functions}

The dissertation introduces stronger differentiability notions because pointwise difference
quotients do not sufficiently coordinate behavior across infinitesimal scales. A derivate-type
condition controls a two-point remainder uniformly. Schematically, for an intended derivative
$g$ one requires
\[
f(y)-f(x)-g(x)(y-x)
\]
to be small relative to $y-x$ in a way uniform over an appropriate neighborhood or scale range.
The exact regularity conditions specify where the generalized series converge and ensure that
termwise algebra, differentiation, and substitution are legitimate.

An expandable function is locally represented by a convergent generalized Taylor expansion with
controlled coefficients and remainder. The important point is not the label but the closure
package: the class is designed to be stable under algebraic operations and differentiation while
supporting Rolle, Lagrange mean value, extremum, and Taylor-type theorems in corrected form.

``Transcendental function'' is not exclusively a complex-analysis term. It generally means a
function not algebraic over the basic rational-function structure; exponential, logarithmic, and
trigonometric functions are standard examples. Complex analysis studies them powerfully, but they
also exist on real and non-Archimedean domains. On $\LC$, power-series definitions require careful
convergence domains, which is why topology and regularity enter before the function theory.

\section{Distributions and genuine functions}

Saying that a generalized delta becomes a function does not erase distribution theory. It means
the enlarged codomain contains infinite values and infinitesimal widths, so a sharply concentrated
profile can be an honest $\LC$-valued function. To compare it with the classical delta, one still
tests its integral action on suitable test functions and then takes a real part or standard part.
The relationship is therefore representational:
\[
\text{non-Archimedean function}
\xrightarrow{\text{pairing and real-part extraction}}
\text{classical distributional action}.
\]
The benefit is that nonlinear or pointwise algebra may be performed before projecting back to the
classical distributional level. The cost is dependence on the enlarged field, its integration
theory, and the chosen projection.

\section{Computer functions and higher derivatives}

A floating-point number is a finite machine representation, usually of the form
\[
\pm m\,b^e,
\]
with bounded precision for the significand $m$ and bounded exponent $e$. Most real numbers cannot
be represented exactly, and arithmetic is rounded. This explains cancellation in finite
differences and why taking a smaller ordinary step does not indefinitely improve accuracy.

An extension algebra enriches each scalar by formal components that propagate derivative
information. Dual numbers $a+b\varepsilon$ with $\varepsilon^2=0$ are the simplest example:
\[
f(x+\varepsilon)=f(x)+f'(x)\varepsilon.
\]
Truncated Taylor algebras keep higher powers and therefore higher derivatives. Levi--Civita
arithmetic is related but not identical: its infinitesimal is nonzero and not generally nilpotent,
and its generalized series can encode a hierarchy of orders.

For a genuine infinitesimal $d$, a forward difference of order $n$ is
\[
\Delta_d^n f(x_0)=
d^{-n}\sum_{j=0}^{n}(-1)^{n-j}\binom nj f(x_0+jd).
\]
In floating-point arithmetic this expression is vulnerable to rounding and cancellation. In a
formal infinitesimal algebra, error orders occupy different powers of $d$ and can be read
coefficient by coefficient. Higher derivatives matter in accelerator and beam physics because
high-order transfer maps describe nonlinear aberrations, stability, and sensitivity. Chapter 7
therefore realizes the earlier algebraic and analytic theory as a computational method; it is not
an unrelated appendix.

\section{What survives from classical analysis?}

When moving to a new scalar field, the reusable procedure is:
\begin{enumerate}
\item identify the algebraic statement and its field assumptions;
\item identify the topology used in convergence or continuity;
\item locate hidden uses of connectedness, compactness, or completeness;
\item test whether the function class is closed under the required operations;
\item strengthen hypotheses only where a proof mechanism actually fails.
\end{enumerate}

Algebraic identities and finite polynomial manipulations usually transfer readily. Results based
on connected intervals, Heine--Borel compactness, or Dedekind completeness require special care.
Power-series arguments may transfer under weak convergence but not under order convergence, or
vice versa. Mean-value and extremum theorems return only after the domain and function class supply
the missing global control.

\section{Common misconceptions, corrected}

\begin{description}[style=nextline,leftmargin=2.3em]
\item[``Real analysis is being abandoned.'']
No. Its limit theory remains valid. The dissertation enlarges the scalar field for a different
kind of representation and then audits which theorems transfer.
\item[``Left-finite means finite support.'']
No. It means finitely many support points below each fixed exponent bound.
\item[``The leading exponent $3/2$ is derived from a universal formula.'']
No. In the displayed example it is simply the smallest exponent whose coefficient is nonzero.
\item[``Weak convergence is vague convergence.'']
No. It is precise convergence in the topology generated by a specified family of seminorms.
\item[``A point particle proves physical infinitesimals exist.'']
No. Point support, quantized values, and non-Archimedean numbers are different concepts.
\item[``A delta represented by a function makes distributions obsolete.'']
No. The non-Archimedean model supplies an alternative representative whose classical action must
still be related to distribution theory.
\item[``Topological differentiability automatically gives the mean value theorem.'']
No. The classical proof also depends on order, connectedness, compactness, and a sufficiently
coherent function class.
\end{description}

\section{A practical second-reading plan}

Read Chapter 1 for the separation of mathematical, physical, and computational motives. In
Chapter 2, first master Archimedean equivalence classes and value groups; treat automorphisms as a
structural classification problem rather than a list of formulas. Chapter 3 is the foundation:
be able to state left-finiteness, convolution multiplication, leading order, the field order, and
the weak seminorms from memory. In Chapter 4, record explicitly which topology is used in every
convergence theorem. In Chapter 5, ask which familiar real-analysis proof fails and at exactly
which hypothesis. Read Chapter 6 as the repair: expandable functions are engineered to restore a
stable calculus. Finally, read Chapter 7 as the computational payoff of the preceding structure.

The shortest conceptual summary is:
\[
\begin{gathered}
\text{abandon Archimedeanness, not rigor;}\\
\text{restrict support, not all infinite series;}\\
\text{choose topology and regularity to match the theorem;}\\
\text{extract finite information from infinitesimal expansions.}
\end{gathered}
\]

\section*{Primary source}
\addcontentsline{toc}{section}{Primary source}
Khodr M. Shamseddine, \textit{New Elements of Analysis on the Levi-Civita Field},
Ph.D. dissertation, Michigan State University, 1999. Author-hosted copy:
\url{https://www.physics.umanitoba.ca/~khodr/dissert.pdf}.

\end{document}
